% The Petit-Ami remote display protocol paper.
% Build: pdflatex remote.tex (twice, for the table of contents), with
% catalog.tex generated from include/graph_remote.h alongside.
\documentclass[11pt]{article}
\usepackage[margin=1in]{geometry}
\usepackage{longtable}
\usepackage{array}
\setlength{\parindent}{0pt}\setlength{\parskip}{6pt plus 2pt}

\title{The Petit-Ami Remote Display Protocol}
\author{S. A. Franco}
\date{August 2026}

\begin{document}
\maketitle

\begin{abstract}
Remote display mode separates a Petit-Ami program from its display: the
program runs on one machine, the client side, and its windows, drawing,
widgets and input devices live on another, the server side. This paper
defines the wire protocol between the two, carried on the message channels
of the network library. Every call of the graphical API has a message; a
continuous event stream flows back. The protocol is defined by the header
\texttt{include/graph\_remote.h}, of which this paper is the narrative
companion; where the two disagree, the header governs.
\end{abstract}

\tableofcontents

\section{Introduction}

The terminal and graphics libraries have always reserved a place for
``remote display'': the output of a program encoded and carried over a
communications channel, with input events carried back, so that the full
functionality of the API survives the trip. The manual states the two
conditions that make this possible: all output must be applied in the order
issued, and all input events must be received in the order generated, with
the rule that outstanding output completes before input is taken.

Two modules implement the mode:

\begin{description}
\item[\texttt{graph\_client}] links with the program. It exposes the
complete graphical API; each call composes a message and sends it, and
\texttt{event()} returns the next inbound event from the client's queue,
fed by a receiver thread.
\item[\texttt{graph\_server}] owns the display. It plugs into the standard
graphics package as an overrider, executes inbound call messages against
it, and sends the resulting events back out.
\end{description}

The display, the mouse, the keyboard and the joysticks face the user on the
server side, while the program logic runs on the client side. Relative to
the usual naming, where the machine in front of the user is the client,
this arrangement is upside down: the ``server'' is the screen on the desk,
and the ``client'' may be a machine in another room entirely. The naming
here follows the connection roles: the server waits for a connection, the
client makes one.

\section{Architecture}

The link is two message channels on adjacent ports, plus a third port
reserved for file transfer:

\begin{center}
\begin{tabular}{lll}
port $p$   & command channel & client $\rightarrow$ server calls, and
                               server $\rightarrow$ client replies \\
port $p+1$ & event channel   & server $\rightarrow$ client events,
                               continuously \\
port $p+2$ & file port       & stream connection, opened per transfer \\
\end{tabular}
\end{center}

The default command port is 4901.

The server divides its work between two threads. The main thread runs the
command loop: wait for a message, execute it against the display library,
reply if the call returns values. A second thread runs the event pump: call
\texttt{event()} forever and send each event out on the event channel. The
two directions are asynchronous, which is exactly why they are separate
threads; neither ever waits on the other.

The client runs one thread beside the program: a receiver, which takes
event messages from the event channel as they arrive and appends them to
the client's event queue. Calls stay on the program's thread: they go out
on the command channel as they are made, in order, and a query sends its
request and waits for the reply before anything else is issued.
\texttt{event()} takes from the queue, waiting if it is empty.

\section{Transport}

The protocol rides the message channels of the network library:
\texttt{openmsg()} on the client side, \texttt{waitmsg()} on the server
side, \texttt{rdmsg()} and \texttt{wrmsg()} to move discrete messages. One
protocol message is one channel message; nothing is split across datagrams,
and nothing shares one.

Message channels bound the size of a message, found from \texttt{maxmsg()}
for the address. The write stream chunks itself to fit; every other message
must fit the bound with its payload, which in practice bounds string
lengths. Senders are responsible for staying under the bound.

The protocol requires a reliable channel in the sense of
\texttt{relymsg()}: no loss, no duplication, no reordering. The local
machine connection of the standard test arrangement is reliable by
definition. Operation over an unreliable network needs a retransmission
layer beneath this protocol; that layer is future work, and nothing in the
message design precludes it.

The secure flag of the message channel calls passes through: a secured
deployment secures both channels and the file connections alike.

\section{Wire format}

All multibyte values are little endian. The scalar and composite notations
used throughout the catalog:

\begin{center}
\begin{tabular}{ll}
\texttt{i64} & signed 64 bit integer, the API \texttt{long} \\
\texttt{f64} & IEEE 754 double, carrying the API float values \\
\texttt{str} & i32 byte length, then the bytes, no terminator \\
\texttt{blk} & i32 byte length, then the bytes \\
\texttt{slst} & string list: i32 count, then count \texttt{str} \\
\texttt{menu} & menu tree, defined below \\
\texttt{evt} & event record, defined below \\
\end{tabular}
\end{center}

Every message begins with a fixed sixteen byte header:

\begin{center}
\begin{tabular}{lll}
\texttt{i32} & \texttt{len} & total message length in bytes, header
included \\
\texttt{i16} & \texttt{mid} & message id, from the catalog \\
\texttt{i16} & \texttt{seq} & request serial, echoed by the reply; 0 where
not meaningful \\
\texttt{i64} & \texttt{wid} & window handle the call acts on; 0 where not
meaningful \\
\end{tabular}
\end{center}

The payload follows, holding the call's parameters in the order the API
declares them, with the window file parameter excluded: it is the header
\texttt{wid}.

\subsection{Menu trees}

A menu serializes depth first: an \texttt{i32} count, then that many
entries, where each entry is an \texttt{i64} flag word (bit 0 the on/off
capability, bit 1 the one-of membership, bit 2 the bar), the \texttt{i64}
id, the \texttt{str} face text, and then a nested menu for the branch,
empty if the entry has none. An empty tree is a single zero count, and
removes the menu where that has meaning.

\texttt{stdmenu()} is the one call that returns a menu: the client sends
the standard selections and the program's additions, and the server
returns the combined tree as it built it. The client materializes that
tree as real menu records for the program, so the program can hand it back
to \texttt{menu()} later, as the API contract requires.

\subsection{Event records}

An event is a fixed record of eleven \texttt{i64}: the window handle the
event belongs to, the event code, the handled flag, and eight parameter
slots holding the fields of the event union in declaration order. A
character rides in the first slot. The slots not used by an event code are
zero. The fixed layout spends a few bytes to buy a single marshaling
routine for every event.

\section{Session}

The first message on the command channel is the hello: the client sends its
protocol version and the server replies with its own. A server that does
not speak the client's version replies with its version and closes; the
client raises the failure to the program. The first message on the event
channel is the event-channel hello from the client, which is how the
server learns where events go; it carries the version again and has no
reply.

The client announces the end of the session with a bye on the command
channel; the server drops the connection. A server error that cannot be
returned as a reply travels to the client as an error message on the event
channel, carrying the error text; the client raises it to the program.

\section{Calls, queries and ordering}

Calls divide in two. A command carries no result: it is sent and the
client proceeds. A query returns values: the client sends the request and
waits on the command channel for the reply before issuing anything else.

There is therefore at most one request outstanding at any time. Replies
cannot interleave with each other or with commands, and the server may
process every message in arrival order with no reassembly. The
\texttt{seq} serial in the header is echoed by the reply purely as a
consistency check; a mismatch is a protocol failure.

The ordinary output of a program, the characters written to a window file,
travels as the write stream message: a block of bytes exactly as they
would have gone to the window, chunked by the sender to the channel bound.
Order is preserved end to end: the write stream and the drawing calls
arrive in the order the program issued them, which satisfies the manual's
first condition for remote display. The event stream preserves generation
order, which satisfies the second.

\section{Windows}

Window handles are small integers assigned by the client. Handle 1 is the
main window, the standard input and output pair, and exists from the
hello. \texttt{openwin()} carries the parent handle, the new handle, and
the logical window id; the server maps each handle to a window file of its
own. Every subsequent call names its window by handle in the message
header. Closing a window file on the client side sends the close message,
and the server closes its side.

\section{Events and client-local operations}

Events are pushed, never polled. A pull model, where the program asks for
the next event and waits for it, would degenerate to an ask and answer
protocol: a full round trip across the wire per event, which for a stream
of mouse movements is a performance disaster. Instead the source runs
ahead of the consumer at both ends. The server's event pump thread calls
\texttt{event()} on the display library without pause, translating each
event to the wire record and sending it the moment it is ready. The
client's receiver thread takes each event message as it arrives and
appends it to the client's event queue, so the transport is always
drained; \texttt{event()} takes from the queue, translates back, and
returns the event to the program, replacing the wire handle with the
window identification the API promises. The program regulates its own
consumption from the queue.

The queue is also where the anticipated optimization lives: redundant
events, a run of mouse movements of which only the last matters, can
collapse in the queue, and a superseding message can override queued ones,
all without any change to the protocol on the wire.

Three operations act on the client's own event stream and put nothing on
the wire: \texttt{eventover()} and \texttt{eventsover()}, which hook event
delivery, hook the client's delivery; and \texttt{sendevent()}, which
injects an event, injects into the client's queue. The input device
queries, the number of mice, buttons, joysticks and function keys, are
synchronous queries like any others; the devices themselves are the
server's, and their activity arrives as events.

\section{File transfer}

Files named in calls live where the program lives, on the client; the
picture files of \texttt{loadpict()} are the case in point. The server
holds a cache, its current directory. When a call names a file, the server
looks it up: first the name as given, then the base name in the cache. If
it holds the file, the call proceeds.

If not, the server requests the file, on the model of ftp's control and
data split. The request goes by message on the command channel, inside the
reply window of the pending call, so the single threaded client is
guaranteed to be listening. The client answers with a port number, then
listens on that port; the server opens the port as a stream connection
with \texttt{opennet()} and receives the file as raw bytes, close
delimited, exactly as an ftp data connection carries one file. The server
stores the file in its cache under its base name, never under a path the
client supplied, and the pending call then proceeds against the cache and
sends its reply.

In this protocol version the port named is always the file port, the
command port plus two; the message names it explicitly so the mechanism
stays general. One transfer at a time falls out of the one outstanding
request rule. The client cannot listen before it sends the port message,
so the server retries the connection briefly to cover the gap.

The exchange is deliberately not tied to pictures: any call that names a
file uses it, and it stands as the protocol's general bulk transfer
mechanism.

\section{Future work}

The protocol as specified assumes a reliable channel and same width, little
endian hosts at both ends, which covers the intended deployments and the
standard test arrangement. Three extensions are anticipated and none
disturb the message catalog: a retransmission layer beneath the protocol
for unreliable networks; print files (\texttt{openprint()}), reserved in
the catalog, whose composition would run on either side; and big endian
hosts, which would swap on the wire boundary. Multiple clients to one
server, and compression of the write stream, remain open questions.

\section{Message catalog}

The catalog below is generated from \texttt{include/graph\_remote.h},
which is the normative definition. Payload notation is as defined in the
wire format section; ``$\rightarrow$'' introduces the reply payload where
the call is a query.

\begin{longtable}{@{}>{\raggedright\arraybackslash}p{2.4in}r>{\raggedright\arraybackslash}p{3.2in}@{}}
\hline
\textbf{Message} & \textbf{Id} & \textbf{Payload $\rightarrow$ reply} \\
\hline
\endhead
% Generated from include/graph_remote.h: the message catalog rows.
% Regenerate with doc/mkremotecatalog.py; the header is normative.
\texttt{GR\_MHELLO} & 1 & first on the command channel: i64 version $\rightarrow$ i64 version. The server refuses a version it does not speak by replying its own and closing. \\
\texttt{GR\_MEVOPEN} & 2 & first on the event channel, from the client, so the server learns where events go: i64 version. No reply. \\
\texttt{GR\_MBYE} & 3 & client is done; the server drops the connection. No payload. \\
\texttt{GR\_MERROR} & 4 & server $\rightarrow$ client on the event channel: str message. The client raises the error to the program. \\
\texttt{GR\_MFILEREQ} & 5 & server $\rightarrow$ client on the command channel, inside the reply window of a pending call that names a file the server does not hold: str name, i64 port, the server's file port. The client opens the port as a stream connection and streams the file in, raw bytes, close delimited; the pending call then completes and replies. \\
\texttt{GR\_MFILEPORT} & 6 & reserved: the active form of the transfer, where the server would connect to a port the client names. The passive form above is the form of this protocol version. \\
\texttt{GR\_MSYNC} & 7 & client $\rightarrow$ server: $\rightarrow$ i64 0. The flow fence: the client bounds its outstanding bytes by syncing once per window, so a command burst cannot outrun the server's socket into datagram loss. Any query is also a fence, being a round trip. \\
\texttt{GR\_MREPLY} & 9 & the reply to any query: the request's seq, and the payload the request's catalog entry gives. \\
\texttt{GR\_MWRITE} & 10 & the write stream of the window: blk of characters, as they would go to the window file. Chunked by the sender to the channel maximum. No reply. \\
\texttt{GR\_MCLOSEWIN} & 11 & the window file closes: the server closes its side. No payload. \\
\texttt{GR\_MCURSOR} & 20 & i64 x, i64 y \\
\texttt{GR\_MMAXX} & 21 & $\rightarrow$ i64 \\
\texttt{GR\_MMAXY} & 22 & $\rightarrow$ i64 \\
\texttt{GR\_MHOME} & 23 & (no payload) \\
\texttt{GR\_MDEL} & 24 & (no payload) \\
\texttt{GR\_MUP} & 25 & (no payload) \\
\texttt{GR\_MDOWN} & 26 & (no payload) \\
\texttt{GR\_MLEFT} & 27 & (no payload) \\
\texttt{GR\_MRIGHT} & 28 & (no payload) \\
\texttt{GR\_MBLINK} & 29 & i64 e \\
\texttt{GR\_MREVERSE} & 30 & i64 e \\
\texttt{GR\_MUNDERLINE} & 31 & i64 e \\
\texttt{GR\_MSUPERSCRIPT} & 32 & i64 e \\
\texttt{GR\_MSUBSCRIPT} & 33 & i64 e \\
\texttt{GR\_MITALIC} & 34 & i64 e \\
\texttt{GR\_MBOLD} & 35 & i64 e \\
\texttt{GR\_MSTRIKEOUT} & 36 & i64 e \\
\texttt{GR\_MSTANDOUT} & 37 & i64 e \\
\texttt{GR\_MFCOLOR} & 38 & i64 color \\
\texttt{GR\_MBCOLOR} & 39 & i64 color \\
\texttt{GR\_MAUTO} & 40 & i64 e \\
\texttt{GR\_MCURVIS} & 41 & i64 e \\
\texttt{GR\_MSCROLL} & 42 & i64 x, i64 y \\
\texttt{GR\_MCURX} & 43 & $\rightarrow$ i64 \\
\texttt{GR\_MCURY} & 44 & $\rightarrow$ i64 \\
\texttt{GR\_MCURBND} & 45 & $\rightarrow$ i64 \\
\texttt{GR\_MSELECT} & 46 & i64 u, i64 d \\
\texttt{GR\_MTIMER} & 47 & i64 i, i64 t, i64 r \\
\texttt{GR\_MKILLTIMER} & 48 & i64 i \\
\texttt{GR\_MMOUSE} & 49 & $\rightarrow$ i64 \\
\texttt{GR\_MMOUSEBUTTON} & 50 & i64 m $\rightarrow$ i64 \\
\texttt{GR\_MJOYSTICK} & 51 & $\rightarrow$ i64 \\
\texttt{GR\_MJOYBUTTON} & 52 & i64 j $\rightarrow$ i64 \\
\texttt{GR\_MJOYAXIS} & 53 & i64 j $\rightarrow$ i64 \\
\texttt{GR\_MSETTAB} & 54 & i64 t \\
\texttt{GR\_MRESTAB} & 55 & i64 t \\
\texttt{GR\_MCLRTAB} & 56 & (no payload) \\
\texttt{GR\_MFUNKEY} & 57 & $\rightarrow$ i64 \\
\texttt{GR\_MFRAMETIMER} & 58 & i64 e \\
\texttt{GR\_MAUTOHOLD} & 59 & i64 e; wid 0, the hold is global \\
\texttt{GR\_MWRTSTR} & 60 & str \\
\texttt{GR\_MWRTSTRN} & 61 & str \\
\texttt{GR\_MSIZBUF} & 62 & i64 x, i64 y \\
\texttt{GR\_MTITLE} & 63 & str \\
\texttt{GR\_MFCOLORC} & 64 & i64 r, i64 g, i64 b \\
\texttt{GR\_MBCOLORC} & 65 & i64 r, i64 g, i64 b \\
\texttt{GR\_MMAXXG} & 100 & $\rightarrow$ i64 \\
\texttt{GR\_MMAXYG} & 101 & $\rightarrow$ i64 \\
\texttt{GR\_MCURXG} & 102 & $\rightarrow$ i64 \\
\texttt{GR\_MCURYG} & 103 & $\rightarrow$ i64 \\
\texttt{GR\_MLINE} & 104 & i64 x1, y1, x2, y2 \\
\texttt{GR\_MLINEWIDTH} & 105 & i64 w \\
\texttt{GR\_MLINESTYLE} & 106 & i64 style \\
\texttt{GR\_MRECT} & 107 & i64 x1, y1, x2, y2 \\
\texttt{GR\_MFRECT} & 108 & i64 x1, y1, x2, y2 \\
\texttt{GR\_MRRECT} & 109 & i64 x1, y1, x2, y2, xs, ys \\
\texttt{GR\_MFRRECT} & 110 & i64 x1, y1, x2, y2, xs, ys \\
\texttt{GR\_MELLIPSE} & 111 & i64 x1, y1, x2, y2 \\
\texttt{GR\_MFELLIPSE} & 112 & i64 x1, y1, x2, y2 \\
\texttt{GR\_MARC} & 113 & i64 x1, y1, x2, y2, sa, ea \\
\texttt{GR\_MFARC} & 114 & i64 x1, y1, x2, y2, sa, ea \\
\texttt{GR\_MFCHORD} & 115 & i64 x1, y1, x2, y2, sa, ea \\
\texttt{GR\_MFTRIANGLE} & 116 & i64 x1, y1, x2, y2, x3, y3 \\
\texttt{GR\_MCURSORG} & 117 & i64 x, i64 y \\
\texttt{GR\_MBASELINE} & 118 & $\rightarrow$ i64 \\
\texttt{GR\_MSETPIXEL} & 119 & i64 x, i64 y \\
\texttt{GR\_MFOVER} & 120 & (no payload) \\
\texttt{GR\_MBOVER} & 121 & (no payload) \\
\texttt{GR\_MFINVIS} & 122 & (no payload) \\
\texttt{GR\_MBINVIS} & 123 & (no payload) \\
\texttt{GR\_MFXOR} & 124 & (no payload) \\
\texttt{GR\_MBXOR} & 125 & (no payload) \\
\texttt{GR\_MFAND} & 126 & (no payload) \\
\texttt{GR\_MBAND} & 127 & (no payload) \\
\texttt{GR\_MFOR} & 128 & (no payload) \\
\texttt{GR\_MBOR} & 129 & (no payload) \\
\texttt{GR\_MCHRSIZX} & 130 & $\rightarrow$ i64 \\
\texttt{GR\_MCHRSIZY} & 131 & $\rightarrow$ i64 \\
\texttt{GR\_MFONTS} & 132 & $\rightarrow$ i64 \\
\texttt{GR\_MFONT} & 133 & i64 fc \\
\texttt{GR\_MFONTNAM} & 134 & i64 fc $\rightarrow$ str \\
\texttt{GR\_MFONTSIZ} & 135 & i64 s \\
\texttt{GR\_MSETPOINTS} & 136 & f64 ps \\
\texttt{GR\_MPOINTS} & 137 & $\rightarrow$ f64 \\
\texttt{GR\_MCHRSPCY} & 138 & i64 s \\
\texttt{GR\_MCHRSPCX} & 139 & i64 s \\
\texttt{GR\_MDPMX} & 140 & $\rightarrow$ i64 \\
\texttt{GR\_MDPMY} & 141 & $\rightarrow$ i64 \\
\texttt{GR\_MSTRSIZ} & 142 & str $\rightarrow$ i64 \\
\texttt{GR\_MCHRPOS} & 143 & str, i64 p $\rightarrow$ i64 \\
\texttt{GR\_MWRITEJUST} & 144 & str, i64 n \\
\texttt{GR\_MJUSTPOS} & 145 & str, i64 p, i64 n $\rightarrow$ i64 \\
\texttt{GR\_MCONDENSED} & 146 & i64 e \\
\texttt{GR\_MEXTENDED} & 147 & i64 e \\
\texttt{GR\_MXLIGHT} & 148 & i64 e \\
\texttt{GR\_MLIGHT} & 149 & i64 e \\
\texttt{GR\_MXBOLD} & 150 & i64 e \\
\texttt{GR\_MHOLLOW} & 151 & i64 e \\
\texttt{GR\_MRAISED} & 152 & i64 e \\
\texttt{GR\_MSETTABG} & 153 & i64 t \\
\texttt{GR\_MRESTABG} & 154 & i64 t \\
\texttt{GR\_MFCOLORG} & 155 & i64 r, i64 g, i64 b \\
\texttt{GR\_MBCOLORG} & 156 & i64 r, i64 g, i64 b \\
\texttt{GR\_MLOADPICT} & 157 & i64 p, str name $\rightarrow$ i64 0. The server fills its cache by the file transfer exchange if it does not hold the file; the reply comes after the picture is loaded. \\
\texttt{GR\_MPICTSIZX} & 158 & i64 p $\rightarrow$ i64 \\
\texttt{GR\_MPICTSIZY} & 159 & i64 p $\rightarrow$ i64 \\
\texttt{GR\_MPICTURE} & 160 & i64 p, x1, y1, x2, y2 \\
\texttt{GR\_MDELPICT} & 161 & i64 p \\
\texttt{GR\_MSCROLLG} & 162 & i64 x, i64 y \\
\texttt{GR\_MPATH} & 163 & i64 a \\
\texttt{GR\_MVIEWOFFG} & 164 & i64 x, i64 y \\
\texttt{GR\_MVIEWSCALE} & 165 & f64 x, f64 y \\
\texttt{GR\_MSCALEX} & 166 & i64 x $\rightarrow$ i64 \\
\texttt{GR\_MSCALEY} & 167 & i64 y $\rightarrow$ i64 \\
\texttt{GR\_MBLOCKCOPYG} & 168 & i64 s, d, sx1, sy1, sx2, sy2, dx1, dy1, dx2, dy2 \\
\texttt{GR\_MOPENWIN} & 200 & open a window: i64 parent handle (0 for none), i64 new handle, i64 logical window id. The header wid is 0. \\
\texttt{GR\_MBUFFER} & 201 & i64 e \\
\texttt{GR\_MSIZBUFG} & 202 & i64 x, i64 y \\
\texttt{GR\_MGETSIZ} & 203 & $\rightarrow$ i64 x, i64 y \\
\texttt{GR\_MGETSIZG} & 204 & $\rightarrow$ i64 x, i64 y \\
\texttt{GR\_MSETSIZ} & 205 & i64 x, i64 y \\
\texttt{GR\_MSETSIZG} & 206 & i64 x, i64 y \\
\texttt{GR\_MSETPOS} & 207 & i64 x, i64 y \\
\texttt{GR\_MSETPOSG} & 208 & i64 x, i64 y \\
\texttt{GR\_MDRAGWIN} & 209 & (no payload) \\
\texttt{GR\_MSCNSIZ} & 210 & $\rightarrow$ i64 x, i64 y \\
\texttt{GR\_MSCNSIZG} & 211 & $\rightarrow$ i64 x, i64 y \\
\texttt{GR\_MSCNCEN} & 212 & $\rightarrow$ i64 x, i64 y \\
\texttt{GR\_MSCNCENG} & 213 & $\rightarrow$ i64 x, i64 y \\
\texttt{GR\_MWINCLIENT} & 214 & i64 cx, cy, ms $\rightarrow$ i64 wx, wy \\
\texttt{GR\_MWINCLIENTG} & 215 & i64 cx, cy, ms $\rightarrow$ i64 wx, wy \\
\texttt{GR\_MFRONT} & 216 & (no payload) \\
\texttt{GR\_MBACK} & 217 & (no payload) \\
\texttt{GR\_MFRAME} & 218 & i64 e \\
\texttt{GR\_MSIZABLE} & 219 & i64 e \\
\texttt{GR\_MSYSBAR} & 220 & i64 e \\
\texttt{GR\_MMENU} & 221 & menu; an empty tree removes the menu \\
\texttt{GR\_MMENUENA} & 222 & i64 id, i64 e \\
\texttt{GR\_MMENUSEL} & 223 & i64 id, i64 e \\
\texttt{GR\_MSTDMENU} & 224 & i64 standard menu set, menu the user additions $\rightarrow$ menu, the combined menu as the server built it, which the client materializes for the program \\
\texttt{GR\_MGETWINID} & 225 & $\rightarrow$ i64; wid 0, ids are global \\
\texttt{GR\_MFOCUS} & 226 & (no payload) \\
\texttt{GR\_MGETWIGID} & 300 & $\rightarrow$ i64 \\
\texttt{GR\_MKILLWIDGET} & 301 & i64 id \\
\texttt{GR\_MSELECTWIDGET} & 302 & i64 id, i64 e \\
\texttt{GR\_MENABLEWIDGET} & 303 & i64 id, i64 e \\
\texttt{GR\_MGETWIDGETTEXT} & 304 & i64 id $\rightarrow$ str \\
\texttt{GR\_MPUTWIDGETTEXT} & 305 & i64 id, str \\
\texttt{GR\_MSIZWIDGET} & 306 & i64 id, i64 x, i64 y \\
\texttt{GR\_MSIZWIDGETG} & 307 & i64 id, i64 x, i64 y \\
\texttt{GR\_MPOSWIDGET} & 308 & i64 id, i64 x, i64 y \\
\texttt{GR\_MPOSWIDGETG} & 309 & i64 id, i64 x, i64 y \\
\texttt{GR\_MBACKWIDGET} & 310 & i64 id \\
\texttt{GR\_MFRONTWIDGET} & 311 & i64 id \\
\texttt{GR\_MFOCUSWIDGET} & 312 & i64 id \\
\texttt{GR\_MBUTTONSIZ} & 313 & str $\rightarrow$ i64 w, i64 h \\
\texttt{GR\_MBUTTONSIZG} & 314 & str $\rightarrow$ i64 w, i64 h \\
\texttt{GR\_MBUTTON} & 315 & i64 x1, y1, x2, y2, str, i64 id \\
\texttt{GR\_MBUTTONG} & 316 & i64 x1, y1, x2, y2, str, i64 id \\
\texttt{GR\_MCHECKBOXSIZ} & 317 & str $\rightarrow$ i64 w, i64 h \\
\texttt{GR\_MCHECKBOXSIZG} & 318 & str $\rightarrow$ i64 w, i64 h \\
\texttt{GR\_MCHECKBOX} & 319 & i64 x1, y1, x2, y2, str, i64 id \\
\texttt{GR\_MCHECKBOXG} & 320 & i64 x1, y1, x2, y2, str, i64 id \\
\texttt{GR\_MRADIOBUTTONSIZ} & 321 & str $\rightarrow$ i64 w, i64 h \\
\texttt{GR\_MRADIOBUTTONSIZG} & 322 & str $\rightarrow$ i64 w, i64 h \\
\texttt{GR\_MRADIOBUTTON} & 323 & i64 x1, y1, x2, y2, str, i64 id \\
\texttt{GR\_MRADIOBUTTONG} & 324 & i64 x1, y1, x2, y2, str, i64 id \\
\texttt{GR\_MGROUPSIZG} & 325 & str, i64 cw, i64 ch $\rightarrow$ i64 w, h, ox, oy \\
\texttt{GR\_MGROUPSIZ} & 326 & str, i64 cw, i64 ch $\rightarrow$ i64 w, h, ox, oy \\
\texttt{GR\_MGROUP} & 327 & i64 x1, y1, x2, y2, str, i64 id \\
\texttt{GR\_MGROUPG} & 328 & i64 x1, y1, x2, y2, str, i64 id \\
\texttt{GR\_MBACKGROUND} & 329 & i64 x1, y1, x2, y2, i64 id \\
\texttt{GR\_MBACKGROUNDG} & 330 & i64 x1, y1, x2, y2, i64 id \\
\texttt{GR\_MSCROLLVERTSIZG} & 331 & $\rightarrow$ i64 w, i64 h \\
\texttt{GR\_MSCROLLVERTSIZ} & 332 & $\rightarrow$ i64 w, i64 h \\
\texttt{GR\_MSCROLLVERT} & 333 & i64 x1, y1, x2, y2, i64 id \\
\texttt{GR\_MSCROLLVERTG} & 334 & i64 x1, y1, x2, y2, i64 id \\
\texttt{GR\_MSCROLLHORIZSIZG} & 335 & $\rightarrow$ i64 w, i64 h \\
\texttt{GR\_MSCROLLHORIZSIZ} & 336 & $\rightarrow$ i64 w, i64 h \\
\texttt{GR\_MSCROLLHORIZ} & 337 & i64 x1, y1, x2, y2, i64 id \\
\texttt{GR\_MSCROLLHORIZG} & 338 & i64 x1, y1, x2, y2, i64 id \\
\texttt{GR\_MSCROLLPOS} & 339 & i64 id, i64 r \\
\texttt{GR\_MSCROLLSIZ} & 340 & i64 id, i64 r \\
\texttt{GR\_MNUMSELBOXSIZG} & 341 & i64 l, i64 u $\rightarrow$ i64 w, i64 h \\
\texttt{GR\_MNUMSELBOXSIZ} & 342 & i64 l, i64 u $\rightarrow$ i64 w, i64 h \\
\texttt{GR\_MNUMSELBOX} & 343 & i64 x1, y1, x2, y2, l, u, id \\
\texttt{GR\_MNUMSELBOXG} & 344 & i64 x1, y1, x2, y2, l, u, id \\
\texttt{GR\_MEDITBOXSIZG} & 345 & str $\rightarrow$ i64 w, i64 h \\
\texttt{GR\_MEDITBOXSIZ} & 346 & str $\rightarrow$ i64 w, i64 h \\
\texttt{GR\_MEDITBOX} & 347 & i64 x1, y1, x2, y2, i64 id \\
\texttt{GR\_MEDITBOXG} & 348 & i64 x1, y1, x2, y2, i64 id \\
\texttt{GR\_MPROGBARSIZG} & 349 & $\rightarrow$ i64 w, i64 h \\
\texttt{GR\_MPROGBARSIZ} & 350 & $\rightarrow$ i64 w, i64 h \\
\texttt{GR\_MPROGBAR} & 351 & i64 x1, y1, x2, y2, i64 id \\
\texttt{GR\_MPROGBARG} & 352 & i64 x1, y1, x2, y2, i64 id \\
\texttt{GR\_MPROGBARPOS} & 353 & i64 id, i64 pos \\
\texttt{GR\_MLISTBOXSIZG} & 354 & slst $\rightarrow$ i64 w, i64 h \\
\texttt{GR\_MLISTBOXSIZ} & 355 & slst $\rightarrow$ i64 w, i64 h \\
\texttt{GR\_MLISTBOX} & 356 & i64 x1, y1, x2, y2, slst, i64 id \\
\texttt{GR\_MLISTBOXG} & 357 & i64 x1, y1, x2, y2, slst, i64 id \\
\texttt{GR\_MDROPBOXSIZG} & 358 & slst $\rightarrow$ i64 cw, ch, ow, oh \\
\texttt{GR\_MDROPBOXSIZ} & 359 & slst $\rightarrow$ i64 cw, ch, ow, oh \\
\texttt{GR\_MDROPBOX} & 360 & i64 x1, y1, x2, y2, slst, i64 id \\
\texttt{GR\_MDROPBOXG} & 361 & i64 x1, y1, x2, y2, slst, i64 id \\
\texttt{GR\_MDROPEDITBOXSIZG} & 362 & slst $\rightarrow$ i64 cw, ch, ow, oh \\
\texttt{GR\_MDROPEDITBOXSIZ} & 363 & slst $\rightarrow$ i64 cw, ch, ow, oh \\
\texttt{GR\_MDROPEDITBOX} & 364 & i64 x1, y1, x2, y2, slst, i64 id \\
\texttt{GR\_MDROPEDITBOXG} & 365 & i64 x1, y1, x2, y2, slst, i64 id \\
\texttt{GR\_MSLIDEHORIZSIZG} & 366 & $\rightarrow$ i64 w, i64 h \\
\texttt{GR\_MSLIDEHORIZSIZ} & 367 & $\rightarrow$ i64 w, i64 h \\
\texttt{GR\_MSLIDEHORIZ} & 368 & i64 x1, y1, x2, y2, i64 mark, i64 id \\
\texttt{GR\_MSLIDEHORIZG} & 369 & i64 x1, y1, x2, y2, i64 mark, i64 id \\
\texttt{GR\_MSLIDEVERTSIZG} & 370 & $\rightarrow$ i64 w, i64 h \\
\texttt{GR\_MSLIDEVERTSIZ} & 371 & $\rightarrow$ i64 w, i64 h \\
\texttt{GR\_MSLIDEVERT} & 372 & i64 x1, y1, x2, y2, i64 mark, i64 id \\
\texttt{GR\_MSLIDEVERTG} & 373 & i64 x1, y1, x2, y2, i64 mark, i64 id \\
\texttt{GR\_MTABBARSIZG} & 374 & slst, i64 tor, i64 cw, i64 ch $\rightarrow$ i64 w, h, ox, oy \\
\texttt{GR\_MTABBARSIZ} & 375 & slst, i64 tor, i64 cw, i64 ch $\rightarrow$ i64 w, h, ox, oy \\
\texttt{GR\_MTABBARCLIENTG} & 376 & i64 tor, i64 w, i64 h $\rightarrow$ i64 cw, ch, ox, oy \\
\texttt{GR\_MTABBARCLIENT} & 377 & i64 tor, i64 w, i64 h $\rightarrow$ i64 cw, ch, ox, oy \\
\texttt{GR\_MTABBAR} & 378 & i64 x1, y1, x2, y2, slst, i64 tor, i64 id \\
\texttt{GR\_MTABBARG} & 379 & i64 x1, y1, x2, y2, slst, i64 tor, i64 id \\
\texttt{GR\_MTABSEL} & 380 & i64 id, i64 tn \\
\texttt{GR\_MALERT} & 420 & str title, str message $\rightarrow$ nothing; the reply comes when the dialog is dismissed, which is what makes the call modal at the far end \\
\texttt{GR\_MQUERYCOLOR} & 421 & i64 r, g, b $\rightarrow$ i64 r, g, b; wid 0 \\
\texttt{GR\_MQUERYOPEN} & 422 & str $\rightarrow$ str; wid 0 \\
\texttt{GR\_MQUERYSAVE} & 423 & str $\rightarrow$ str; wid 0 \\
\texttt{GR\_MQUERYFIND} & 424 & str, i64 opt $\rightarrow$ str, i64 opt; wid 0 \\
\texttt{GR\_MQUERYFINDREP} & 425 & str, str, i64 opt $\rightarrow$ str, str, i64 opt; wid 0 \\
\texttt{GR\_MQUERYFONT} & 426 & i64 fc, s, fr, fg, fb, br, bg, bb, effect $\rightarrow$ the same nine, as chosen \\
\texttt{GR\_MEVENT} & 500 & an event, server $\rightarrow$ client on the event channel: evt \\
\hline
\multicolumn{3}{@{}l}{\textbf{The sound partition (base 8192)}} \\
\hline
\texttt{GR\_MSTARTTIMEOUT} & 8192 & -- \\
\texttt{GR\_MSTOPTIMEOUT} & 8193 & -- \\
\texttt{GR\_MCURTIMEOUT} & 8194 & $\rightarrow$ i64 time \\
\texttt{GR\_MSTARTTIMEIN} & 8195 & -- \\
\texttt{GR\_MSTOPTIMEIN} & 8196 & -- \\
\texttt{GR\_MCURTIMEIN} & 8197 & $\rightarrow$ i64 time \\
\texttt{GR\_MSYNTHOUT} & 8198 & $\rightarrow$ i64 count \\
\texttt{GR\_MSYNTHIN} & 8199 & $\rightarrow$ i64 count \\
\texttt{GR\_MOPENSYNTHOUT} & 8200 & i64 p \\
\texttt{GR\_MCLOSESYNTHOUT} & 8201 & i64 p \\
\texttt{GR\_MOPENSYNTHIN} & 8202 & i64 p \\
\texttt{GR\_MCLOSESYNTHIN} & 8203 & i64 p \\
\texttt{GR\_MNOTEON} & 8204 & i64 p, t, c, n, v \\
\texttt{GR\_MNOTEOFF} & 8205 & i64 p, t, c, n, v \\
\texttt{GR\_MINSTCHANGE} & 8206 & i64 p, t, c, i \\
\texttt{GR\_MATTACK} & 8207 & i64 p, t, c, at \\
\texttt{GR\_MRELEASE} & 8208 & i64 p, t, c, rt \\
\texttt{GR\_MLEGATO} & 8209 & i64 p, t, c, b \\
\texttt{GR\_MPORTAMENTO} & 8210 & i64 p, t, c, b \\
\texttt{GR\_MVIBRATO} & 8211 & i64 p, t, c, v \\
\texttt{GR\_MVOLSYNTHCHAN} & 8212 & i64 p, t, c, v \\
\texttt{GR\_MPORTTIME} & 8213 & i64 p, t, c, v \\
\texttt{GR\_MBALANCE} & 8214 & i64 p, t, c, b \\
\texttt{GR\_MPAN} & 8215 & i64 p, t, c, b \\
\texttt{GR\_MTIMBRE} & 8216 & i64 p, t, c, tb \\
\texttt{GR\_MBRIGHTNESS} & 8217 & i64 p, t, c, b \\
\texttt{GR\_MREVERB} & 8218 & i64 p, t, c, r \\
\texttt{GR\_MTREMULO} & 8219 & i64 p, t, c, tr \\
\texttt{GR\_MCHORUS} & 8220 & i64 p, t, c, cr \\
\texttt{GR\_MCELESTE} & 8221 & i64 p, t, c, ce \\
\texttt{GR\_MPHASER} & 8222 & i64 p, t, c, ph \\
\texttt{GR\_MAFTERTOUCH} & 8223 & i64 p, t, c, n, at \\
\texttt{GR\_MPRESSURE} & 8224 & i64 p, t, c, pr \\
\texttt{GR\_MPITCH} & 8225 & i64 p, t, c, pt \\
\texttt{GR\_MPITCHRANGE} & 8226 & i64 p, t, c, v \\
\texttt{GR\_MMONO} & 8227 & i64 p, t, c, ch \\
\texttt{GR\_MPOLY} & 8228 & i64 p, t, c \\
\texttt{GR\_MLOADSYNTH} & 8229 & i64 s, str fn $\rightarrow$ i64 0 \\
\texttt{GR\_MPLAYSYNTH} & 8230 & i64 p, t, s \\
\texttt{GR\_MDELSYNTH} & 8231 & i64 s \\
\texttt{GR\_MWAITSYNTH} & 8232 & i64 p $\rightarrow$ i64 0, on completion \\
\texttt{GR\_MWRSYNTH} & 8233 & i64 p + seq: i64 port, time, st, p1, p2, p3 \\
\texttt{GR\_MRDSYNTH} & 8234 & i64 p $\rightarrow$ seq; blocks the caller until input arrives, as the native call does \\
\texttt{GR\_MWAVEOUT} & 8235 & $\rightarrow$ i64 count \\
\texttt{GR\_MWAVEIN} & 8236 & $\rightarrow$ i64 count \\
\texttt{GR\_MOPENWAVEOUT} & 8237 & i64 p \\
\texttt{GR\_MCLOSEWAVEOUT} & 8238 & i64 p \\
\texttt{GR\_MLOADWAVE} & 8239 & i64 w, str fn $\rightarrow$ i64 0 \\
\texttt{GR\_MPLAYWAVE} & 8240 & i64 p, t, w \\
\texttt{GR\_MDELWAVE} & 8241 & i64 w \\
\texttt{GR\_MVOLWAVE} & 8242 & i64 p, t, v \\
\texttt{GR\_MWAITWAVE} & 8243 & i64 p $\rightarrow$ i64 0, on completion \\
\texttt{GR\_MCHANWAVEOUT} & 8244 & i64 p, c \\
\texttt{GR\_MRATEWAVEOUT} & 8245 & i64 p, r \\
\texttt{GR\_MLENWAVEOUT} & 8246 & i64 p, l \\
\texttt{GR\_MSGNWAVEOUT} & 8247 & i64 p, s \\
\texttt{GR\_MFLTWAVEOUT} & 8248 & i64 p, f \\
\texttt{GR\_MENDWAVEOUT} & 8249 & i64 p, e \\
\texttt{GR\_MWRWAVE} & 8250 & i64 p, i64 frames, blk samples $\rightarrow$ i64 0. The sender chunks to the channel bound on frame boundaries; each chunk is acknowledged, which paces the stream to the consumer and keeps it whole. \\
\texttt{GR\_MOPENWAVEIN} & 8251 & i64 p \\
\texttt{GR\_MCLOSEWAVEIN} & 8252 & i64 p \\
\texttt{GR\_MCHANWAVEIN} & 8253 & i64 p $\rightarrow$ i64 \\
\texttt{GR\_MRATEWAVEIN} & 8254 & i64 p $\rightarrow$ i64 \\
\texttt{GR\_MLENWAVEIN} & 8255 & i64 p $\rightarrow$ i64 \\
\texttt{GR\_MSGNWAVEIN} & 8256 & i64 p $\rightarrow$ i64 \\
\texttt{GR\_MENDWAVEIN} & 8257 & i64 p $\rightarrow$ i64 \\
\texttt{GR\_MFLTWAVEIN} & 8258 & i64 p $\rightarrow$ i64 \\
\texttt{GR\_MRDWAVE} & 8259 & i64 p, frames $\rightarrow$ i64 got, blk samples; the server clamps the answer to the channel bound and the caller loops. Lengths are in samples, one frame of every channel, as the API counts. \\
\texttt{GR\_MSYNTHOUTNAME} & 8260 & i64 p $\rightarrow$ str \\
\texttt{GR\_MSYNTHINNAME} & 8261 & i64 p $\rightarrow$ str \\
\texttt{GR\_MWAVEOUTNAME} & 8262 & i64 p $\rightarrow$ str \\
\texttt{GR\_MWAVEINNAME} & 8263 & i64 p $\rightarrow$ str \\
\texttt{GR\_MSETPARAMSYNTHIN} & 8264 & i64 p, str, str $\rightarrow$ i64 \\
\texttt{GR\_MSETPARAMSYNTHOUT} & 8265 & i64 p, str, str $\rightarrow$ i64 \\
\texttt{GR\_MSETPARAMWAVEIN} & 8266 & i64 p, str, str $\rightarrow$ i64 \\
\texttt{GR\_MSETPARAMWAVEOUT} & 8267 & i64 p, str, str $\rightarrow$ i64 \\
\texttt{GR\_MGETPARAMSYNTHIN} & 8268 & i64 p, str $\rightarrow$ str \\
\texttt{GR\_MGETPARAMSYNTHOUT} & 8269 & i64 p, str $\rightarrow$ str \\
\texttt{GR\_MGETPARAMWAVEIN} & 8270 & i64 p, str $\rightarrow$ str \\
\texttt{GR\_MGETPARAMWAVEOUT} & 8271 & i64 p, str $\rightarrow$ str \\

\hline
\end{longtable}

\end{document}
